% ============================================================================
% Cabeçalho LaTeX do PDF consolidado do ciclo (Regra 16).
% Este arquivo é um MODELO: os marcadores @NOME@, @CICLO@ e @PERIODO@ são
% substituídos por tools/gerar-ciclo.sh antes de chamar o pandoc.
%
% Objetivo: rodapé com identificação do ciclo e número de página em TODAS as
% páginas (inclusive a capa) e índice começando em página nova.
%
% Observação de ambiente (ver README-ciclos.md): nesta máquina o engine usado é
% o pdflatex, porque o xelatex não está instalado e o lualatex está quebrado
% (falta lualatex-math.sty e o luaotfload erra ao carregar as fontes OpenType).
% Por isso aqui não se usa fontspec: apenas recursos clássicos, que renderizam
% acentuação portuguesa corretamente com o pacote babel-portuguese.
% ============================================================================

\usepackage{fancyhdr}

% Texto do rodapé: identificação do ciclo + número da página
\newcommand{\rodapeciclo}{@NOME@ \textperiodcentered\ p.\ \thepage}

% Estilo aplicado ao documento inteiro
\fancypagestyle{ciclo}{%
  \fancyhf{}%
  \fancyfoot[C]{\footnotesize\rodapeciclo}%
  \renewcommand{\headrulewidth}{0pt}%
  \renewcommand{\footrulewidth}{0pt}%
}

% A capa (primeira página) usa o estilo 'plain' por padrão: redefinimos para
% que ela também receba o rodapé, como pede a regra de entrega.
\fancypagestyle{plain}{%
  \fancyhf{}%
  \fancyfoot[C]{\footnotesize\rodapeciclo}%
  \renewcommand{\headrulewidth}{0pt}%
  \renewcommand{\footrulewidth}{0pt}%
}

\pagestyle{ciclo}

% O índice abre em página própria (capa na p. 1, índice na p. 2)
\let\cicloTocOriginal\tableofcontents
\renewcommand{\tableofcontents}{\clearpage\cicloTocOriginal}

% Título do índice e profundidade
\renewcommand{\contentsname}{Índice}
\setcounter{tocdepth}{3}

% Espaçamento um pouco mais generoso entre parágrafos do que o padrão LaTeX
\setlength{\parskip}{0.4em}

% ---------------------------------------------------------------------------
% Compatibilidade Unicode: símbolos que o pdflatex não traz por padrão.
% Sem isto, o PDF falha com "Unicode character ... not set up for use with LaTeX"
% (foi o caso de ≠, ✓, ●, ↺ e ⌒ usados nos documentos do ciclo).
% ---------------------------------------------------------------------------
\usepackage{amssymb}
\usepackage{newunicodechar}
\newunicodechar{≠}{\ensuremath{\neq}}
\newunicodechar{✓}{\ensuremath{\checkmark}}
\newunicodechar{●}{\textbullet}
\newunicodechar{↺}{\ensuremath{\circlearrowleft}}
\newunicodechar{⌒}{\ensuremath{\frown}}
